%\documentclass[12pt, letterpaper]{book} % Para una tesis book es buuena opción
%\documentclass[12pt, letterpaper]{report} %Para una tesis report es buena opción
\documentclass[12pt,a4paper]{article} % Configuración básica: tamaño de fuente y papel

%% =========================
%% PAQUETES ESENCIALES
%% =========================
\usepackage[spanish, es-tabla, es-nodecimaldot]{babel} % Soporte para español 
\usepackage[utf8]{inputenc}    % Codificación de caracteres UTF-8 (pdflatex)
\usepackage[T1]{fontenc}       % Mejor renderizado de fuentes
\usepackage{lmodern}           % Usa la fuente Latin Modern
\usepackage{microtype}         % Mejoras tipográficas sutiles
\usepackage{coffeestains}      % Marca de café jaja

%% =========================
%% MATEMÁTICAS Y FÍSICA
%% =========================
\usepackage{amsmath}           % Entornos matemáticos avanzados y símbolos
\usepackage{amssymb}           % Símbolos matemáticos adicionales
\usepackage{amsfonts}          % Fuentes matemáticas adicionales (como ℝ, ℤ, etc.)
\usepackage{mathtools}         % Extensiones y correcciones para amsmath
\usepackage{physics}           % Notación física (derivadas, vectores, etc.)
\usepackage{siunitx}           % Formateo correcto de unidades (S.I.)
\sisetup{per-mode = fraction}  % Configuración de unidades como fracciones

%% ==========================
%% QUÍMICA
%% ==========================
\usepackage[version=4]{mhchem}

%% =========================
%% GRÁFICOS Y TABLAS
%% =========================
\usepackage{graphicx}         % Inclusión de imágenes
    \graphicspath{{figs/}}    % (buena práctica: carpeta para figuras)
\usepackage{array}            % Más opciones para tablas y matrices
\usepackage{booktabs}         % Mejores tablas (formato profesional)
\usepackage{caption}          % Personalización de leyendas
\usepackage{subcaption}       % Sub-figuras y sub-tablas
\usepackage{tabularx}         % tablas con ancho flexible (útil)
\usepackage{longtable}        % tablas que ocupan varias páginas (si aplica)
\usepackage{adjustbox}        % Manipulación avanzada de cajas y contenido
\usepackage{tcolorbox}        % cajas estilizadas (notas, ejemplos, resultados)

\usepackage{tikz}             % dibujo vectorial y diagramas
%\usetikzlibrary{shapes.geometric, arrows.meta, positioning, babel}
\usepackage{pgfplots}         % gráficos (espectros, curvas)
\pgfplotsset{compat=1.18}

%% =========================
%% REFERENCIAS Y METADATOS
%% =========================
\usepackage{hyperref}         % Hipervínculos en el documento
\hypersetup{colorlinks=true, linkcolor=midnightblue, citecolor=red,
            pdftitle={Título del artículo}, pdfauthor={Aguilar A.}}

\usepackage{csquotes}         % recomendado con biblatex para citas correctas
\usepackage[spanish]{cleveref}         % referencias inteligentes (\cref, \Cref)


%% =========================
%% FORMATO Y TIPOGRAFÍA
%% =========================
\usepackage{xcolor}
\usepackage{geometry}         % Márgenes personalizados
\geometry{
    margin=2.0cm,   % Configuración de márgenes
    headheight=15pt, % Aumentamos la altura del encabezado
    includehead      % Opcional: incluye el encabezado en el cálculo del margen
         }
\usepackage{fancyhdr}
\pagestyle{fancy}
\fancyhf{}
\fancyhead[L]{\small Aguilar A.}
\fancyhead[R]{\small Título del Artículo}
\fancyfoot[C]{\thepage}

%% =========================
%% OTROS PAQUETES ÚTILES
%% =========================
\usepackage{float}            % Mejor control de posicionamiento de flotantes
\usepackage{lipsum}           % Texto de relleno (¡ELIMINAR en documento final!)
\usepackage{lineno}           % numerar líneas en borradores si lo deseas

%% =========================
%% PAQUETES ADICIONALES
%% =========================

%% =========================
%% Datos del Paper
%% =========================
\title{Título del artículo}
\author{Aguilar A. \quad y \quad Coautor \\
        \small{Centro de Investigaciones Químicas} \\
        \small{Universidad Autónoma del Estado de Morelos}}
\date{\today}

% ----------------------------------------------------------------------------
\begin{document}

\maketitle
% Añade una mancha de café en la primera página
% {opacidad}{escala}{ángulo}{offset-x}{offset-y}
\coffeestainA{0.8}{0.75}{-45}{5cm}{10.5cm}  

%\tableofcontents

\begin{abstract}
    \lipsum[1] % Texto de ejemplo. Reemplazar con resumen real
\end{abstract}

\section{Introducción}
\lipsum[2] % Texto de ejemplo

\section{Marco Teórico}
\lipsum[3]

\section{Manejo de unidades}
Unidades SI: \SI{9.8}{\meter\per\second\squared}

\section{Metodología}

\section{Resultados y Discusión}
\begin{table}[H] % begin cuando queremos manipular la tabla como [H]
    \centering
    \caption{Parámetros del estudio}
    \begin{tabular}{@{}llr@{}}
        %{@{}llr@{}} % quitar espacios antes y después del texto
        \toprule
        Material & Formula & Densidad (\unit{\gram\per\centi\meter\cubed}) \\
        \midrule
        Agua     & Agua      & 1.0 \\
        Hierro   & \ce{Fe}   & 7.87 \\
        \bottomrule
    \end{tabular}
\end{table}

\begin{tabular}{|l|c|r|} % Las barras | crean líneas verticales
   \hline
   \textbf{Producto} & \textbf{Cantidad} & \textbf{Precio} \\
   \hline
   Manzanas & 10 & 20.50 \\
   \hline
   Peras & 5 & 15.00 \\
   \hline
\end{tabular}


% Bibliografía
\bibliographystyle{achemso} % estilo de bibliografía (se puede borrar)
\bibliography{bibliografia}

\end{document}

%% =========================
%% Información para textos
%% =========================

%% Graficos
%    \centering                                                      
%         \includegraphics[width=0.8\linewidth]{ejemplo.png}
%    \caption{Ejemplo de figura}
%
%% Cuerpo del documento en dos columnas
%    \begin{multicols}{2}
%    \end{multicols}
%

%% Citas - Leer ~/tesis/tutorial/latex/bibliografia/bib_resume/readme.txt

%%  Contenido de bibliografia.bib
% @article{einstein1905,
%   author  = {Albert Einstein},
%   title   = {Zur Elektrodynamik bewegter Körper},
%   journal = {Annalen der Physik},
%   volume  = {322},
%   number  = {10},  <-- No la incluyó Minhhuy en un ejp.
%   pages   = {891--921},
%   year    = {1905},
% }

%% Herramientas
